\documentclass[12pt]{article}
\usepackage[utf8]{inputenc}
\usepackage[T1]{fontenc}
\usepackage{hyperref}
\title{Kamu Language Resources for K-12 Curriculum Development}
\author{Compiled by Claude Sonnet 5.5 (claude-sonnet-5-5) for Andr\'as Kornai}
\date{}
\begin{document}
\maketitle

\section{Introduction}
Kamu (ISO 639-3 xmu; Glottocode kamu1258), also called Gamor, Kamor or Malak Malak, is an Indigenous Australian language of the Northern Territory, classified in Eastern Daly. Wikipedia reports two speakers in 1975 and extinction by the 1980s. Glottolog classes it as extinct, while the Endangered Languages Catalogue still lists it as critically endangered. Some authors treat it as a dialect of Matngala. My searches found no curriculum materials, primers or audio for Kamu.

\section{Primary Educational Resources}
No primers, textbooks or school programmes for Kamu were found in search results.

\section{Linguistic and Reference Materials}
\begin{itemize}
\item Wikipedia, ``Kamu language'': \url{https://en.wikipedia.org/wiki/Kamu_language}
\item Wikipedia, ``Eastern Daly languages'': \url{https://en.wikipedia.org/wiki/Eastern_Daly_languages}
\item Glottolog entry for Kamu: \url{https://glottolog.org/resource/languoid/id/kamu1258}
\item Endangered Languages Catalogue entry, with a bibliography page: \url{https://endangeredlanguages.com/lang/3239?hl=en} and \url{https://endangeredlanguages.com/lang/3239/bibliography}
\item ASJP word list for Kamor: \url{https://asjp.clld.org/languages/KAMOR}
\item Mark Harvey (2003) bibliographic records in Glottolog (the search result did not show the titles): \url{https://glottolog.org/resource/reference/id/ozbib:2658} and \url{https://glottolog.org/resource/reference/id/ozbib:2659}
\end{itemize}

\section{Religious and Community Resources}
\begin{itemize}
\item Wikipedia, ``Kamu people'': \url{https://en.wikipedia.org/wiki/Kamu_people}
\end{itemize}
No community organisation or scripture was found.

\section{Digital and Audio Resources}
No audio or video was found. The Verbix page lists the language among endangered languages: \url{https://docs.verbix.com/EndangeredLanguages/Kamu}

\section{Curriculum Development Context}
Very little is documented for Kamu, and it has no known speakers. Any heritage or revitalization curriculum would have to be built from archival word lists and the sources by Harvey and others listed above. The closely related Eastern Daly languages, such as Matngala, may offer comparative material, but I did not find such resources. Descendants should be involved in any such work.
\end{document}